\chapter{Gut Health}

\section*{Definition and Overview}
Gut health refers to the wellbeing of the digestive system, particularly the large intestine (colon) and the trillions of bacteria, viruses, and fungi that live there, collectively called the gut microbiome. A healthy gut digests food efficiently, absorbs nutrients, protects against harmful organisms, and maintains a balanced microbiome that supports immunity and even brain function. Gut health is influenced by diet, exercise, sleep, stress, medicines, and illness. Symptoms such as bloating, irregular bowel habits, and discomfort are common and are often improved by dietary and lifestyle changes. There is growing evidence that gut health affects overall health, though "gut health" products and advice vary widely in quality.

\section*{Epidemiology}
Digestive symptoms are among the most common reasons for GP visits. Around 1 in 5 Australians experiences symptoms of irritable bowel syndrome, and functional gut symptoms such as bloating and irregular bowel habits affect up to a third of the population. The gut microbiome is a rapidly growing area of research, and interest in gut health has increased enormously, with probiotics, prebiotics, and fibre supplements among the most popular health products. Conditions linked to gut health include inflammatory bowel disease, irritable bowel syndrome, obesity, and some immune and mental health conditions.

\section*{Aetiology and Pathophysiology}
Gut health reflects the complex interaction of the gut, its microbiome, and the rest of the body.

\begin{itemize}
    \item \textbf{The microbiome:} trillions of bacteria in the colon ferment fibre, produce vitamins and short-chain fatty acids, and interact with the immune system
    \item \textbf{The gut barrier:} the lining of the gut keeps harmful substances out while absorbing nutrients
    \item \textbf{Gut-brain axis:} the gut and brain communicate through nerves, hormones, and immune signals, linking gut health to mood and stress
    \item \textbf{Diet:} fibre feeds beneficial bacteria; highly processed foods and excess sugar reduce microbial diversity
    \item \textbf{Medicines:} antibiotics can disrupt the microbiome, and some medicines cause gut symptoms
    \item \textbf{Lifestyle:} sleep, exercise, and stress affect gut function
    \item \textbf{Disease:} infection, inflammation, and food intolerances affect gut health
\end{itemize}

\section*{Clinical Features}
Signs that suggest gut health may be suboptimal include:

\begin{itemize}
    \item Bloating, wind, and a feeling of fullness after meals
    \item Irregular bowel habits, including constipation, diarrhoea, or alternating
    \item Abdominal discomfort or cramping
    \item Changes in stool consistency or frequency
    \item Heartburn and indigestion
    \item Food intolerances or sensitivities
    \item Fatigue, which can be linked to poor digestion and absorption
    \item Frequent infections, which some research links to immune and gut imbalance
    \item A sudden change in bowel habit or blood in the stool, which requires medical assessment
\end{itemize}

\section*{Differential Diagnosis}
Gut symptoms can reflect conditions that need specific treatment.

\begin{itemize}
    \item Irritable bowel syndrome, a common functional disorder
    \item Inflammatory bowel disease, including Crohn disease and ulcerative colitis
    \item Coeliac disease and other food intolerances
    \item Gastroenteritis and intestinal infections
    \item Gastro-oesophageal reflux disease
    \item Gallbladder and pancreatic disease
    \item Bowel cancer, which must be considered with bleeding, weight loss, or a change in bowel habit in older people
    \item Medicine-related gut symptoms
    \item Thyroid and other metabolic conditions affecting the gut
\end{itemize}

\section*{When to Seek Medical Care}
See a doctor for persistent or severe gut symptoms, sudden changes in bowel habit, blood in the stool, unexplained weight loss, or symptoms that wake you from sleep. These can indicate conditions requiring investigation.

\textbf{Call 000 (or local emergency services) for:}
\begin{itemize}
    \item Severe abdominal pain, especially with vomiting or a rigid abdomen
    \item Vomiting blood or passing black, tarry stools
    \item Severe dehydration from vomiting or diarrhoea
    \item High fever with abdominal pain
    \item Signs of bowel obstruction, including an inability to pass gas or stool with vomiting and distension
\end{itemize}

\section*{Diagnosis and Investigations}
Assessment depends on the symptoms and any alarm features.

\begin{itemize}
    \item \textbf{History and examination:} symptom pattern, diet, medicines, and family history
    \item \textbf{Blood tests:} full blood count, coeliac antibodies, inflammatory markers, and thyroid function
    \item \textbf{Stool tests:} for infection, inflammation, or blood
    \item \textbf{Breath tests:} for lactose intolerance and bacterial overgrowth
    \item \textbf{Endoscopy and colonoscopy:} for alarm symptoms or suspected inflammatory bowel disease
    \item \textbf{Imaging:} ultrasound, CT, or MRI where appropriate
    \item \textbf{Food and symptom diary:} helps identify triggers
\end{itemize}

\section*{Management}
Improving gut health centres on diet and lifestyle, with medical treatment for underlying conditions.

\begin{itemize}
    \item \textbf{Eat more fibre:} a variety of fruits, vegetables, legumes, and whole grains feeds beneficial bacteria
    \item \textbf{Prebiotics:} foods such as onions, garlic, bananas, and oats support good bacteria
    \item \textbf{Probiotics:} fermented foods such as yoghurt and kefir, or supplements, can support the microbiome in some people
    \item \textbf{Limit processed foods:} reduce highly processed foods, excess sugar, and artificial sweeteners
    \item \textbf{Stay hydrated:} adequate water supports regular bowel function
    \item \textbf{Exercise regularly:} activity stimulates normal gut motility
    \item \textbf{Manage stress:} stress affects the gut through the gut-brain axis
    \item \textbf{Sleep well:} consistent sleep supports gut and overall health
    \item \textbf{Medical treatment:} management of underlying conditions, including irritable bowel syndrome and inflammatory bowel disease
    \item \textbf{Review medicines:} discuss any medicines affecting the gut with your doctor
\end{itemize}

\section*{Complications}
\begin{itemize}
    \item Malnutrition and deficiencies when absorption is impaired
    \item Dehydration from chronic diarrhoea
    \item Progression of undiagnosed conditions such as inflammatory bowel disease or cancer
    \item Reduced quality of life from chronic symptoms
    \item Wasted spending on unproven "gut health" products
    \item Effects of repeated antibiotic courses on the microbiome
    \item Increased risk of complications from untreated disease
\end{itemize}

\section*{Prognosis and Outcomes}
For most people, gut health improves with simple dietary and lifestyle changes, and symptoms such as bloating and irregular bowel habits respond well to fibre, hydration, exercise, and stress management. People with specific conditions such as irritable bowel syndrome or coeliac disease usually improve markedly with appropriate treatment. The gut microbiome is adaptable, and changes in diet can improve microbial diversity within weeks. The key is a balanced approach, professional assessment of persistent symptoms, and caution with unproven products.

\section*{Prevention and Self-Care}
\begin{itemize}
    \item Eat a varied diet rich in fibre from plant foods
    \item Include fermented foods and probiotics regularly
    \item Drink plenty of water and limit alcohol
    \item Exercise most days
    \item Manage stress with relaxation, mindfulness, or talking therapies
    \item Prioritise sleep and a consistent routine
    \item Eat regular meals and chew food well
    \item Avoid unnecessary antibiotics, and take them only as prescribed
    \item Keep a food and symptom diary if symptoms are troublesome
    \item Seek medical advice for persistent, severe, or worrying symptoms
\end{itemize}

\section*{Sources and Further Reading}
The following reputable sources provide further information for healthcare students and patients seeking reliable, current advice about gut health.

\begin{itemize}
    \item Gastroenterological Society of Australia. \textit{Digestive health information for patients.} https://www.gesa.org.au/
    \item Dietitians Australia. \textit{Gut health and the microbiome.} https://dietitiansaustralia.org.au/
    \item Healthdirect Australia. \textit{Gut health.} https://www.healthdirect.gov.au/
    \item CSIRO. \textit{Gut health research and resources.} https://www.csiro.au/
    \item NHS. \textit{Good foods to help your digestion.} https://www.nhs.uk/
    \item Harvard T.H. Chan School of Public Health. \textit{The gut microbiome.} https://www.hsph.harvard.edu/
\end{itemize}
